\documentclass{article}
\usepackage{lmodern}
\usepackage{amsmath}
\usepackage{listings}
\usepackage{fullpage}
\usepackage{parskip}
\usepackage{xcolor}
\lstset{
	basicstyle=\ttfamily,
	columns=fullflexible,
	frame=single,
	breaklines=true,
	postbreak=\mbox{\textcolor{red}{$\hookrightarrow$}\space},
}
\begin{document}
\title{Experiment `p\_6\_2\_e\_remove\_middle\_seq' Results}
\author{\today}
\date{}
\maketitle
\textbf{Experiment outcome: }SUCCESS
\\\textbf{Bad responses: }0
\\\textbf{Responses containing}~\texttt{assume}~\textbf{: }0
\\\textbf{Responses containing}~\texttt{expect}~\textbf{: }0
\\\textbf{Resolution attempts: }2
\\\textbf{Hard fails (resolution): }0
\\\textbf{Soft fails (resolution): }0
\\\textbf{Verification attempts: }2
\section*{Problem Specification}
\textbf{Problem name: }p\_6\_2\_e\_remove\_middle\_seq
\\\textbf{Natural language statement: }Write a method to remove the middle element of an array if the array length is odd, or the middle two elements if the length is even.
\\\textbf{Method signature: }p\_6\_2\_e\_remove\_middle\_seq(inputs: seq<int>) returns (result: seq<int>)
\subsection*{Ensures}
\begin{itemize}
\item \texttt{|result| == if |inputs| \% 2 == 1 then |inputs| - 1 else |inputs| - 2}
\item \texttt{|inputs| \% 2 == 1 ==> forall i :: 0 <= i < |inputs| / 2 ==> result[i] == inputs[i] \&\& forall j :: |inputs| / 2 <= j < |result| ==> result[j] == inputs[j + 1]}
\item \texttt{|inputs| \% 2 == 0 ==> forall i :: 0 <= i < |inputs| / 2 - 1 ==> result[i] == inputs[i] \&\& forall j :: |inputs| / 2 - 1 <= j < |result| ==> result[j] == inputs[j + 2]}
\end{itemize}
\subsection*{Requires}
\begin{itemize}
\item \texttt{|inputs| >= 1}
\end{itemize}
\clearpage
\section*{GenAI interactions}
Below you will find all interactions between the `user' (program) and the `assistant' (OpenAI).
\subsection*{Program $\rightarrow$ GenAI}
\begin{lstlisting}
You are given the following task to perform in Dafny:

Write a method to remove the middle element of an array if the array length is odd, or the middle two elements if the length is even.

The signature should be:

p_6_2_e_remove_middle_seq(inputs: seq<int>) returns (result: seq<int>)

The method should respect the following contract:

requires |inputs| >= 1, ensures |result| == if |inputs| % 2 == 1 then |inputs| - 1 else |inputs| - 2, ensures |inputs| % 2 == 1 ==> forall i :: 0 <= i < |inputs| / 2 ==> result[i] == inputs[i] && forall j :: |inputs| / 2 <= j < |result| ==> result[j] == inputs[j + 1], ensures |inputs| % 2 == 0 ==> forall i :: 0 <= i < |inputs| / 2 - 1 ==> result[i] == inputs[i] && forall j :: |inputs| / 2 - 1 <= j < |result| ==> result[j] == inputs[j + 2]

Produce and show only the Dafny body of this method, including the curly braces that surround it. Do not show the signature nor contract. You must not use 'assume' nor 'expect' anywhere in your code.
\end{lstlisting}
\subsection*{GenAI $\rightarrow$ Program}
\textbf{System fingerprint: }null
\\\textbf{ID: }chatcmpl-Dn0HMVXwONBbW8K0vMGOO6nHlYcpQ
\\\textbf{Model: }gpt-5.5-2026-04-23
\\\textbf{Created at: }1780571516
\\\textbf{Finish reason: }stop
\\\textbf{Usage: }606 tokens in, and 1186 tokens out
\begin{lstlisting}
{
  if |inputs| % 2 == 1 {
    var m := |inputs| / 2;
    var newLen := |inputs| - 1;
    result := seq(newLen, k => if k < m then inputs[k] else inputs[k + 1]);
  } else {
    assert |inputs| >= 2;
    var m := |inputs| / 2;
    var newLen := |inputs| - 2;
    result := seq(newLen, k => if k < m - 1 then inputs[k] else inputs[k + 2]);
  }
}
\end{lstlisting}
\subsection*{Program $\rightarrow$ GenAI}
\begin{lstlisting}
When using dafny verify, the below error is emitted and verify fails:

problem_data/p_6_2_e_remove_middle_seq/program.dfy(10,51): Error: index out of range
   |
10 |     result := seq(newLen, k => if k < m then inputs[k] else inputs[k + 1]);
   |                                                    ^

problem_data/p_6_2_e_remove_middle_seq/program.dfy(10,66): Error: index out of range
   |
10 |     result := seq(newLen, k => if k < m then inputs[k] else inputs[k + 1]);
   |                                                                   ^

problem_data/p_6_2_e_remove_middle_seq/program.dfy(15,55): Error: index out of range
   |
15 |     result := seq(newLen, k => if k < m - 1 then inputs[k] else inputs[k + 2]);
   |                                                        ^

problem_data/p_6_2_e_remove_middle_seq/program.dfy(15,70): Error: index out of range
   |
15 |     result := seq(newLen, k => if k < m - 1 then inputs[k] else inputs[k + 2]);
   |                                                                       ^


Dafny program verifier finished with 1 verified, 4 errors

Correct the error in code by altering only the method body. Produce and show only the Dafny body, including the curly braces that surround it. Do not show the signature nor contract. You must not use 'assume' nor 'expect' anywhere in your code.
\end{lstlisting}
\subsection*{GenAI $\rightarrow$ Program}
\textbf{System fingerprint: }null
\\\textbf{ID: }chatcmpl-Dn0IEEY0on9StDifo05BeQwhaIC1A
\\\textbf{Model: }gpt-5.5-2026-04-23
\\\textbf{Created at: }1780571570
\\\textbf{Finish reason: }stop
\\\textbf{Usage: }1130 tokens in, and 7021 tokens out
\begin{lstlisting}
{
  var n := |inputs|;
  if n % 2 == 1 {
    var m := n / 2;
    assert n == (n / 2) * 2 + n % 2;
    assert n == m * 2 + 1;
    assert m * 2 == 2 * m;
    assert n == 2 * m + 1;
    assert 0 <= m;
    assert m + 1 <= n;

    var left := inputs[..m];
    var right := inputs[m + 1..];
    result := left + right;

    assert |left| == m;
    assert |right| == n - (m + 1);
    assert |result| == |left| + |right|;
    assert |result| == n - 1;

    forall i | 0 <= i < m
      ensures result[i] == inputs[i]
    {
      assert i < |left|;
      assert result[i] == left[i];
      assert left[i] == inputs[i];
    }

    forall j | m <= j < |result|
      ensures result[j] == inputs[j + 1]
    {
      assert |left| == m;
      assert j >= |left|;
      assert 0 <= j - m;
      assert j < n - 1;
      assert j - m < n - (m + 1);
      assert j - m < |right|;
      assert j - |left| == j - m;
      assert result[j] == right[j - |left|];
      assert result[j] == right[j - m];
      assert right[j - m] == inputs[(m + 1) + (j - m)];
      assert (m + 1) + (j - m) == j + 1;
    }
  } else {
    var m := n / 2;
    assert 0 <= n % 2 < 2;
    assert n % 2 == 0;
    assert n == (n / 2) * 2 + n % 2;
    assert n == m * 2;
    assert m * 2 == 2 * m;
    assert n == 2 * m;
    assert 0 <= m;
    if m < 1 {
      assert m == 0;
      assert n == 0;
      assert false;
    }
    assert m >= 1;
    assert m - 1 >= 0;
    assert m + 1 <= n;

    var left := inputs[..(m - 1)];
    var right := inputs[m + 1..];
    result := left + right;

    assert |left| == m - 1;
    assert |right| == n - (m + 1);
    assert |result| == |left| + |right|;
    assert |result| == n - 2;

    forall i | 0 <= i < m - 1
      ensures result[i] == inputs[i]
    {
      assert i < |left|;
      assert result[i] == left[i];
      assert left[i] == inputs[i];
    }

    forall j | m - 1 <= j < |result|
      ensures result[j] == inputs[j + 2]
    {
      assert |left| == m - 1;
      assert j >= |left|;
      assert 0 <= j - (m - 1);
      assert j < n - 2;
      assert j - (m - 1) < n - (m + 1);
      assert j - (m - 1) < |right|;
      assert j - |left| == j - (m - 1);
      assert result[j] == right[j - |left|];
      assert result[j] == right[j - (m - 1)];
      assert right[j - (m - 1)] == inputs[(m + 1) + (j - (m - 1))];
      assert (m + 1) + (j - (m - 1)) == j + 2;
    }
  }
}
\end{lstlisting}
\section*{Final Program}
\begin{lstlisting}
method p_6_2_e_remove_middle_seq(inputs: seq<int>) returns (result: seq<int>)
	requires |inputs| >= 1
	ensures |result| == if |inputs| % 2 == 1 then |inputs| - 1 else |inputs| - 2
	ensures |inputs| % 2 == 1 ==> forall i :: 0 <= i < |inputs| / 2 ==> result[i] == inputs[i] && forall j :: |inputs| / 2 <= j < |result| ==> result[j] == inputs[j + 1]
	ensures |inputs| % 2 == 0 ==> forall i :: 0 <= i < |inputs| / 2 - 1 ==> result[i] == inputs[i] && forall j :: |inputs| / 2 - 1 <= j < |result| ==> result[j] == inputs[j + 2]
{
  var n := |inputs|;
  if n % 2 == 1 {
    var m := n / 2;
    assert n == (n / 2) * 2 + n % 2;
    assert n == m * 2 + 1;
    assert m * 2 == 2 * m;
    assert n == 2 * m + 1;
    assert 0 <= m;
    assert m + 1 <= n;

    var left := inputs[..m];
    var right := inputs[m + 1..];
    result := left + right;

    assert |left| == m;
    assert |right| == n - (m + 1);
    assert |result| == |left| + |right|;
    assert |result| == n - 1;

    forall i | 0 <= i < m
      ensures result[i] == inputs[i]
    {
      assert i < |left|;
      assert result[i] == left[i];
      assert left[i] == inputs[i];
    }

    forall j | m <= j < |result|
      ensures result[j] == inputs[j + 1]
    {
      assert |left| == m;
      assert j >= |left|;
      assert 0 <= j - m;
      assert j < n - 1;
      assert j - m < n - (m + 1);
      assert j - m < |right|;
      assert j - |left| == j - m;
      assert result[j] == right[j - |left|];
      assert result[j] == right[j - m];
      assert right[j - m] == inputs[(m + 1) + (j - m)];
      assert (m + 1) + (j - m) == j + 1;
    }
  } else {
    var m := n / 2;
    assert 0 <= n % 2 < 2;
    assert n % 2 == 0;
    assert n == (n / 2) * 2 + n % 2;
    assert n == m * 2;
    assert m * 2 == 2 * m;
    assert n == 2 * m;
    assert 0 <= m;
    if m < 1 {
      assert m == 0;
      assert n == 0;
      assert false;
    }
    assert m >= 1;
    assert m - 1 >= 0;
    assert m + 1 <= n;

    var left := inputs[..(m - 1)];
    var right := inputs[m + 1..];
    result := left + right;

    assert |left| == m - 1;
    assert |right| == n - (m + 1);
    assert |result| == |left| + |right|;
    assert |result| == n - 2;

    forall i | 0 <= i < m - 1
      ensures result[i] == inputs[i]
    {
      assert i < |left|;
      assert result[i] == left[i];
      assert left[i] == inputs[i];
    }

    forall j | m - 1 <= j < |result|
      ensures result[j] == inputs[j + 2]
    {
      assert |left| == m - 1;
      assert j >= |left|;
      assert 0 <= j - (m - 1);
      assert j < n - 2;
      assert j - (m - 1) < n - (m + 1);
      assert j - (m - 1) < |right|;
      assert j - |left| == j - (m - 1);
      assert result[j] == right[j - |left|];
      assert result[j] == right[j - (m - 1)];
      assert right[j - (m - 1)] == inputs[(m + 1) + (j - (m - 1))];
      assert (m + 1) + (j - (m - 1)) == j + 2;
    }
  }
}
\end{lstlisting}
\section*{Total Token Usage}
\textbf{Input tokens: }1736
\\\textbf{Output tokens: }8207
\\\textbf{Reasoning tokens: }6989
\\\textbf{Sum of `total tokens': }9943
\section*{Experiment Timings}
\textbf{Overall Experiment} started at 1780571516459, ended at 1780571725372, lasting 208913ms (208.91 seconds)
\\\textbf{Iteration \#1} started at 1780571516461, ended at 1780571570640, lasting 54179ms (54.18 seconds)
\\\textbf{Iteration \#2} started at 1780571570650, ended at 1780571725372, lasting 154722ms (154.72 seconds)
\end{document}
